\section{Related Work}
\label{sec:related}
\paragraph{Benchmarks and learned localization.}
VCDB and VCSL provide partial-copy localization benchmarks~\cite{jiang2014,vcsl2022};
VCSL scores segment overlap on both axes, but interval labels do not give dense frame
correspondence. VSC 2023 and FiGVCL use different settings or annotations, so their numbers
cannot be merged~\cite{pizzi2024vsc,figvcl2025}. SPD learns a detector on similarity
maps~\cite{spd2021}; TransVCL uses self-/cross-attention, correlation and a localization
head~\cite{transvcl2023}. Its reference inference processes chunk pairs of at most 1,200 rows
per axis; pair-level timing must include all chunks and assembly. RTR trains regional
features and a detector with self-supervised pairs~\cite{rtr2024}; self-supervision does not
make the localizer training-free. Recent variants include iESTA and
VCSA~\cite{ding2025iesta,li2025vcsa}; feature requirements are reported rather than assumed
compatible.

\paragraph{Geometry, alignment and efficient search.}
The VCSL comparison includes HV, TN, DP and DTW~\cite{vcsl2022,tan2009,chou2015}; temporal
consistency and cascading search also appear in retrieval~\cite{douze2016,tian2013}, and SAM
uses temporal line fitting in a matching pipeline~\cite{liu2023sam}. Our zero-floor principle
follows local alignment~\cite{smith1981} with a restricted step set and tube domain; it is not
a new dynamic-programming principle. Temporal-consistency re-ranking~\cite{tcsvt2024} and
library-scale retrieval-plus-matching~\cite{wang2026mltdedup} are adjacent systems. Sampling
efficiency has been studied~\cite{chen2025sampling}; reducing the sampling rate is not
sufficient novelty.

\paragraph{Comparison that matters.}
The proposed distinction is the controlled relationship between sparse candidate support,
finite refinement domains and measured cost, isolated against a simple pipeline with the same
classical ingredients. No priority or performance advantage follows from this review.
